\documentclass[11pt]{article}

\usepackage{geometry}
\usepackage{hyperref}
\usepackage{tikz}
\usetikzlibrary{arrows.meta, positioning, calc}
\geometry{margin=1in}

\title{Emergent State Machines:\\A Clinician's Introduction}
\author{}
\date{}

\begin{document}

\maketitle

\section{The Problem: Too Many Signals, Too Little Structure}

Modern clinical environments generate enormous amounts of patient data. Vital signs, laboratory values, imaging results, and clinical notes are continuously produced and monitored.

Despite this wealth of information, many monitoring systems operate in a very simple way: a signal crosses a threshold, and an alarm is triggered.

Clinicians know the limitations of this approach. Many alerts are not actionable. Important patterns may be buried among less relevant signals. Alarm fatigue is a well-known consequence.

The underlying issue is not a lack of data. It is a lack of structure in how signals are interpreted. See Figure~\ref{fig:clinician-esm-flow}.

\begin{figure}[ht]
\centering
\setlength{\fboxsep}{8pt}
\renewcommand{\arraystretch}{1.35}
\begin{tabular}{c}
\fbox{\parbox{0.76\textwidth}{\centering \textbf{Patient Observations}\\Vital signs, labs, clinical notes, bedside observations}}\\
$\downarrow$\\
\fbox{\parbox{0.76\textwidth}{\centering \textbf{Clinical Signals}\\Direct and derived indicators such as oxygen saturation, respiratory rate trend, shock index, or interaction patterns}}\\
$\downarrow$\\
\fbox{\parbox{0.76\textwidth}{\centering \textbf{Patient State}\\An interpretable representation of the patient's current situation}}\\
$\downarrow$\\
\fbox{\parbox{0.76\textwidth}{\centering \textbf{Clinical Gate}\\Condition determining whether further rule evaluation or escalation should activate}}\\
$\downarrow$\\
\fbox{\parbox{0.76\textwidth}{\centering \textbf{Clinical Policy Projection}\\State expressed in clinically meaningful decision dimensions such as instability, urgency, or escalation priority}}\\
$\downarrow$\\
\fbox{\parbox{0.76\textwidth}{\centering \textbf{Clinical Response}\\Increase monitoring, review history, order diagnostics, notify specialist, or escalate care}}\\
$\downarrow$\\
\fbox{\parbox{0.76\textwidth}{\centering \textbf{Patient Evolution}\\The patient changes over time, producing new observations and a new state}}\\[0.6em]
\parbox{0.82\textwidth}{\centering \textbf{Visible Clinical Trajectory:} clinicians can inspect how the patient's state changes over time and whether interventions alter its direction}
\end{tabular}
\caption{Clinician-oriented ESM flow. The system organizes bedside observations into interpretable patient state and supports clinicians by making evolving situations visible and auditable.}
\label{fig:clinician-esm-flow}
\end{figure}

\section{Situations Instead of Signals}

An Emergent State Machine (ESM) is designed to monitor \textit{situations}, not just individual signals.

Instead of evaluating one measurement at a time, the system evaluates combinations of signals that together represent clinically meaningful conditions.

For example, clinicians often reason about patterns such as:

\begin{itemize}
\item decreasing oxygen saturation
\item increasing respiratory rate
\item rising heart rate
\end{itemize}

Taken together, these signals may indicate worsening respiratory instability.

An ESM represents such patterns explicitly and monitors how they evolve over time.

\section{How the System Works}

The ESM architecture follows a simple reasoning chain:

\begin{itemize}
\item patient observations are recorded
\item observations produce signals
\item signals form a structured patient state
\item clinical policy evaluates the state
\item appropriate actions or alerts are triggered
\end{itemize}

Each step in this process is transparent and inspectable.

If the system produces an alert, clinicians can trace that alert back to the signals and observations that produced it.

\section{Clinical Knowledge Remains in Human Hands}

A key principle of the ESM architecture is that clinical expertise remains central.

Clinicians and medical experts define:

\begin{itemize}
\item which signals matter
\item how signals combine to form meaningful patterns
\item what actions should occur when patterns appear
\end{itemize}

The system performs structured monitoring, but it does not replace clinical judgment.

Instead, it supports clinicians by making relevant situations easier to detect and interpret.

\section{Guided Responses and Clinical Workflows}

When the system detects a developing situation, it can recommend possible next steps.

Examples might include:

\begin{itemize}
\item increasing monitoring frequency
\item reviewing patient history
\item ordering additional tests
\item consulting a specialist
\end{itemize}

These recommendations appear as structured options rather than opaque automated decisions.

Clinicians remain responsible for deciding how to respond.

\section{Why This Matters}

The goal of the ESM approach is not to automate clinical decision-making.

The goal is to support clinicians by:

\begin{itemize}
\item organizing complex signal patterns
\item reducing unnecessary alarms
\item highlighting meaningful clinical situations
\item making system reasoning transparent
\end{itemize}

In environments where clinicians must interpret many signals at once, structured situational monitoring may improve both clarity and workflow.

\section{Looking Forward}

Emergent State Machines are designed as a transparent and governable architecture for complex monitoring systems.

Because every signal, rule, and decision remains inspectable, the system can evolve alongside clinical knowledge while preserving accountability and safety.

The intent is simple: provide better situational awareness while keeping clinical expertise at the center of care.

\end{document}
